\begin{papercard}{arkalakis2024}
\cf{Stages} \stagetag{D} (longitudinal capture of Data Safety sections and of app behaviour), \stagetag{C} (a manual mapping of Data Safety labels to API calls by two coders), \stagetag{I} (their agreement is reported as raw percentages), \stagetag{H} (undisclosed collection and sharing, ignored consent choices). The paper is a USENIX Security 2024 measurement study of whether Google Play's Data Safety section is truthful: an automated dynamic-analysis framework runs Android apps, records which sensitive API calls and network transmissions occur, and compares them with the label the developer declared. It is the Android counterpart of the iOS label studies in this corpus.

\cf{Research problem and contribution} The Data Safety section (in effect since July 2022) relies solely on developer self-declaration, which presumes developers are truthful and fully aware of what their code and third-party libraries do \pc{arkalakis2024}{3}{it relies solely on the developer to declare what data the application will collect and share}. Three contributions: analysis of Google's documentation for ``loopholes''; a mapping of Data Safety labels to function calls; a dynamic pipeline and a longitudinal study of popular apps \pc{arkalakis2024}{3}{Our framework monitors all API methods that lead to sensitive personal information or device data}. Headline: of almost 5K apps, 81\,\% show at least one discrepancy; 69.2\,\% under-declare data collected and 43.1\,\% data shared; 79.4\,\% of apps with incomplete disclosure do not ask for consent and 78.6\,\% of those that ask disregard the choice; 59\,\% of discrepancies come from third-party libraries and 41\,\% from first-party code \pc{arkalakis2024}{3}{69.2\% of apps not fully disclosing the data they collect and 43.1\% the data that they share}.

\cf{Background: structure and loopholes} The section has Data Collected, Data Shared and Security Practices parts and uses predefined labels (Location, Contacts, Photos or videos, \ldots) \pc{arkalakis2024}{3}{The Data Safety section consists of three parts, specifically the Data Collected, Data Shared and Security Practices parts}. Per Google's documentation, ``collected'' means data transmitted \emph{off the device}; exempt are data processed only locally, data end-to-end encrypted, and data processed only ``ephemerally'' \pc{arkalakis2024}{4}{data that is transmitted off device but only processed “ephemerally” does not need to be disclosed}. Sharing exempts transfers to service providers, legal purposes, ``user-initiated action or prominent disclosure and user consent'', and anonymous data. The authors criticise especially the consent exemption, because dark patterns are known in run-time consent dialogs. Definition used: a \textbf{discrepancy} is ``any case where data is accessed or shared at runtime without having been disclosed''; sharing without disclosure is called a \emph{data leak} \pc{arkalakis2024}{4}{Throughout the paper we consider as a discrepancy any case where data is accessed or shared at runtime without having been disclosed in the app’s Data Safety section}.

\cf{Framework} (1) \emph{API monitoring} with Reaper on rooted Pixel devices (Pixel 4a and 6, AOSP 12, Magisk, LSPosed); permission-to-function mappings from DYNAMO (Android 10) extended with the APER method (all @RequiresPermission methods in the Android 12 source) plus known PII functions. (2) \emph{Network monitoring} with mitmproxy, SSL pinning disabled with Objection; Data Shared matched by string search for device-specific values and crafted values (advertising ID, GSF ID, email, name), including common encodings (base64, hex, ROT13, SHA/MD5 variants) \pc{arkalakis2024}{10}{Since the values sent over the network can be in an encoded form, our system checks for common transformations during the matching process}. (3) \emph{UI traversal} with UIHarvester (extended for WebViews), five minutes per app, breadth-first, Google Single Sign-On where possible \pc{arkalakis2024}{6}{two Google Pixel devices (Pixel 4a and Pixel 6) running Google’s AOSP version 12}. (4) A Python \emph{crawler} that logs each app's Data Safety section, permissions, developer and downloads daily from 29 July 2022; the runtime result is compared with the section as crawled on the day the APK was downloaded \pc{arkalakis2024}{6}{We created a Python-based crawler that parses the Google Play web page on a daily basis}. (5) \emph{Consent handling}: a keyword-based detector for runtime consent dialogs, built by manually interacting with 100 apps and checked on another random 100, then five scenarios: no consent dialog; accept; reject; reject button not found; reject and the app exits \pc{arkalakis2024}{6}{We populated this list of keywords by manually interacting with 100 apps that contain runtime consent dialogues}.

\cf{Manual coding: Data Safety labels $\to$ API functions (stage C/I)} Of 6,955 monitored functions, 295 were actually called by the exercised apps; two authors independently assigned Data Safety labels (or NoLabel) using permissions, Google's guideline and the Java source, then compared and resolved differences with a third author by majority \pc{arkalakis2024}{5}{out of the 6,955 function calls that our system monitors, we selected 295 functions that were actually accessed by the apps we exercised}. Result: 94 functions mapped to 17 labels (9 coarse categories, 8 sub-categories); identifiers such as name and email are detected on the network instead. Agreement of the two labellers: \textbf{242 of 295 (82.03\,\%) identical}, 8 (2.71\,\%) partly matching, 44 (14.91\,\%) one label vs.\ NoLabel, 1 true category conflict \pc{arkalakis2024}{5}{had a precise agreement (i.e., all assigned labels match) for 242 functions (82.03\%)}. Only raw percentages are reported, no chance-corrected coefficient. Labels without an API analogue (e.g.\ religious beliefs) cannot be monitored.

\cf{Sample and design} Free apps, up to the top 250 from each of 20 Play categories; 4,986 unique apps, four snapshots (Sep 2022, Mar/Jun/Sep 2023), 21,202 executions in total; Sep 2022 was a preliminary Data Collected-only run; no new apps were added across waves so apps removed from the store shrink later datasets (4,157 and 3,953 apps) \pc{arkalakis2024}{6}{we downloaded 4,986 unique apps from Google Play using the Raccoon [3] framework and performed a total of 21,202 executions}. A \emph{VPN subset} of 100 apps in 10 countries (Canada, USA, Brazil, Iceland, Germany, Ukraine, Estonia, Kenya, Russia, India) re-downloaded the country store version and Data Safety section; geoblocking affected only 18 app-country pairs \pc{arkalakis2024}{6}{We downloaded a subset of 100 apps from our initial dataset and analyzed them in several countries using a VPN service}.

\cf{Results} Discrepancies persisted a year after the section became mandatory: 80.89\,\% (Mar 2023), 80.68\,\% (Jun), 81.46\,\% (Sep) \pc{arkalakis2024}{7}{These discrepancies persist for an entire year after Data Safety was made mandatory}. Fully compliant: 19.1\,\%, 19.3\,\%, 18.5\,\%; entirely non-compliant (declared nothing): about 13--14\,\%; the rest partly declared; mean compliance level of apps with discrepancies 63.9\,\% \pc{arkalakis2024}{7}{794 (19.1\%), 803 (19.32\%) and 729 (18.54\%) apps are fully compliant}. By consent scenario (Sep 2023): 80\,\% of apps with no consent dialog have discrepancies; 84.8\,\% of apps whose user accepts; the reject scenarios show data still collected: of apps with a dialog, 78.61\,\% have discrepancies, and 60--63\,\% of ``reject not found'' apps access data before any user choice \pc{arkalakis2024}{8}{we found that 85.79\% of the apps}. Functions: location, device identifiers (IMEI, MAC, SSID/BSSID, IDs), thermal status; the user's accept/reject decision ``makes almost no difference'' to undisclosed collection. Persistence: only 81 apps fixed their discrepancies between Mar and Jun while 3,282 stayed unchanged and 72 became non-compliant \pc{arkalakis2024}{9}{We found that only 81 apps corrected their discrepancies and are compliant with Data Safety policies}. Data leaks: undisclosed sharing of advertising ID, BSSID, build number, device name, GSF ID, MCC+MNC, SSID and (crafted values) name, email, location to 16,840 domains, mostly ``Internet services'', software/hardware and web-ad categories; 12 apps sent PII over HTTP although their section claimed secure transfer \pc{arkalakis2024}{10}{Interestingly, we found 12 apps sending PIIs to 12 domains over HTTP}. Origin by stack-trace and a third-party-library list: 59\,\% third-party, 41\,\% first-party; 12.34\,\% of apps only first-party and 26.62\,\% only third-party \pc{arkalakis2024}{10}{most of the discrepancies originate from third-party functionality (59\%)}. VPN: 79--86\,\% of apps show discrepancies independent of country (Table 3). Case studies: popular apps with 500M+ installs (TikTok, Instagram, Word/Excel, Telegram, \ldots) have discrepancies; 47 of 99 testable Google apps; of apps with an independent security review only a small number exist, and some still have discrepancies \pc{arkalakis2024}{12}{Out of the remaining 99 Google apps, we found that 47 of them had discrepancies}.

\cf{Recommendations} A universal disclosure standard across Android and iOS, comprehensive policies without exceptions, third-party libraries documenting their data collection, and an automated platform-side validation mechanism with public reports and penalties \pc{arkalakis2024}{13}{Official app markets must not rely solely on the developer’s honesty, and must frequently review app’s run-time behavior}.

\cf{Limits (author and card)} The authors list: the label mapping may be incomplete since Google gives no official mapping \pc{arkalakis2024}{13}{Our mapping between Data Safety labels and API methods may not be complete, as no official documentation exists that maps the labels to function calls}; encrypted or obfuscated traffic is not detected, so Data Shared results are lower bounds \pc{arkalakis2024}{10}{our results for Data Shared discrepancies present a lower bound}; keyword-based consent detection misses cases such as an X-marked decline button; packages missing from the library list default to first-party and obfuscated names inflate first-party counts \pc{arkalakis2024}{13}{package names that are not included in our list will be classified as first parties}. Card additions: five-minute automated traversal exercises only reachable screens (no sign-in beyond Google SSO, no paid features, no mental-health content depth); testing on own rooted devices with test accounts; apps are not fed realistic personal data; discrepancy is defined by Google's own exemptions and the authors' stricter normative view that all data should be disclosed, so some ``violations'' are Google-policy-compliant; the authors do not infer intent; the comparison unit is label vs.\ function, not label vs.\ privacy policy; no chance-corrected reliability for the manual mapping.

\cf{Use for this thesis} (1) The most detailed protocol in the corpus for verifying a store label against behaviour; it shows what a thesis without code-level analysis \emph{cannot} claim: a label or policy can be compared with visible screens, but not with actual data flows. If the thesis keeps to human walkthrough, it should treat the label as a claim, not as ground truth, and cite this paper for the size of the gap (81\,\% discrepancy). (2) The definition of ``discrepancy'' and the distinction between Google's exemptions and a stricter normative standard is a model for stating the thesis' own comparison rule up front. (3) A precedent for double-coding a mapping with a third person resolving differences, but with only percent agreement; the thesis should add a chance-corrected statistic. (4) A model for dated, repeated capture (four waves, daily label crawl, comparison with the label of the download day) and for limiting to a fixed cohort across waves. (5) The three-way comparison label/behaviour/consent reminds the thesis that an app can be consistent with the Data Safety section while still misleading users because of the consent exemption.
\end{papercard}
